\problemname{Circle Merge}
\noindent

Emil has $N$ bags of stones placed in a circle.
Each bag contains some number of stones, where the number of stones in bag $i$ is $a_i$ for $1 \leq i \leq N$.
The bags are arranged in a circle, so that bag $1$ is adjacent to bags $2$ and $N$, bag $2$ is adjacent to bags $1$ and $3$, $\dots$, and bag $N$ is adjacent to bags $N-1$ and $1$.

Emil wants all bags to contain the same number of stones.
In one move, Emil can merge two adjacent bags into a single bag containing the sum of the stones from the two original bags.

What is the minimum number of moves required so that all bags contain the same number of stones?

\section*{Input}
The first line contains an integer $N$ ($1 \le N \le 2 \cdot 10^5$), the number of bags in the circle initially.

The second line contains $N$ integers $a_1, a_2, \ldots, a_N$ ($1 \le a_i \le 5 \cdot 10^{12}$), where $a_i$ is the number of stones in bag $i$.

\section*{Output}
Print a single integer, the minimum number of moves required so that all bags contain the same number of stones.


\section*{Scoring}
Your solution will be tested on a set of test groups, each worth a number of points. Each test group contains
a set of test cases. To get the points for a test group you need to solve all test cases in the test group.

\noindent
\begin{tabular}{| l | l | p{12cm} |}
  \hline
  \textbf{Group} & \textbf{Points} & \textbf{Constraints} \\ \hline
  $1$    & $4$        & There is a solution with the minimum number of moves, where the number of stones in every bag is $a_1$. \\ \hline
  $2$    & $15$       & $N \le 100$ \\ \hline 
  $3$    & $22$       & $N \le 3000, a_i \le 200$ \\ \hline
  $4$    & $29$       & $a_i \le 200$ \\ \hline
  $5$    & $30$       & No additional constraints. \\ \hline 
\end{tabular}


\section*{Explanation of Sample Case 1}
In the first case, it is optimal to merge the bags as follows:
\begin{itemize}
    \item Merge the last and the first bag: [1, 2, 3, 4, 5] $\to$ [6, 2, 3, 4] (note that the bags are arranged in a circle)
    \item Merge the first and second bag: [6, 2, 3, 4] $\to$ [8, 3, 4]
    \item Merge the second and third bag: [8, 3, 4] $\to$ [8, 7]
    \item Merge the last two bags: [8, 7] $\to$ [15].
\end{itemize}

In total, 4 moves are required for all bags to contain the same number of stones. It can be shown that it is not possible to do this in fewer moves.


\section*{Explanation of Sample Case 2}
9 \\
5 2 3 1 4 2 1 1 1

In the second case, it is optimal to merge the bags as follows:
\begin{itemize}
    \item Merge bags 4 and 5: [5, 2, 3, 1, 4, 2, 1, 1, 1] $\to$ [5, 2, 3, 5, 2, 1, 1, 1]
    \item Merge bags 6 and 7: [5, 2, 3, 5, 2, 1, 1, 1] $\to$ [5, 2, 3, 5, 2, 2, 1]
    \item Merge bags 7 and 8: [5, 2, 3, 5, 2, 2, 1] $\to$ [5, 2, 3, 5, 2, 3]
    \item Merge bags 2 and 3: [5, 2, 3, 5, 2, 3] $\to$ [5, 5, 5, 2, 3]
    \item Merge bags 4 and 5: [5, 5, 5, 2, 3] $\to$ [5, 5, 5, 5].
\end{itemize}

In total, 5 moves are required for all bags to contain the same number of stones. It can be shown that it is not possible to do this in fewer moves.

